% =====================================================================
%  Banco complementar --- Capacitores  (REVISADO)
%  Substitui a versao gerada por template (10 clones em N2, 11 em N3,
%  10 questoes conceituais indevidamente marcadas como N4).
%  Sem sobreposicao com o cap10, que ja traz: C de placas paralelas,
%  Q=CV com W, serie 6/3, paralelo 4/12, energia (MC), associacao mista
%  6||3, RC de 1 kohm/1 uF, tres de 6 uF, descarga de 2 uF/6 V, grafico
%  v_C(t), COMPARTILHAMENTO DE CARGA entre C1 e C2 e INSERCAO DE
%  DIELETRICO (a Q e a V constantes) --- estes dois ultimos motivaram a
%  remocao de duas N4 do banco antigo, que os duplicavam.
%  Distribuicao mantida: 4 N1 + 10 N2 + 11 N3 + 10 N4 = 35
% =====================================================================

\subsubsection*{Banco complementar --- Capacitores}

\begin{atencao}[Organização do banco]
Três relações organizam o tópico: $Q=CU$, $W=\frac12CU^{2}$ e
$i=C\dfrac{\mathrm{d}u}{\mathrm{d}t}$. Delas decorre o fato central dos
transitórios --- a \textbf{tensão} no capacitor não pode saltar (isso exigiria
corrente infinita), ao contrário do indutor, em que é a \emph{corrente} que é
contínua. O N3 trata de divisão de tensão em série, fuga, ESR e extração de
parâmetros; o N4 exige dedução da EDO, balanço energético, projeto de banco,
filtro e temporizador.
\end{atencao}

\blocofixacao

\begin{questao}[1]{}
Um capacitor armazena $\SI{300}{\micro\coulomb}$ sob $\SI{20}{\volt}$.
Determine sua capacitância.
\resposta{$C=\SI{15}{\micro\farad}$}
\resolucao{
\[
C=\frac{Q}{U}=\frac{\pot{300}{-6}}{20}=\pot{15}{-6}\,\si{\farad}
=\SI{15}{\micro\farad}.
\]
A capacitância é uma propriedade \emph{geométrica} do componente: ela não
depende de $Q$ nem de $U$, apenas da razão entre eles. Aplicar o dobro da
tensão dobra a carga e deixa $C$ inalterada.}
\end{questao}

\begin{questao}[1]{}
Um capacitor de $\SI{100}{\micro\farad}$ armazena $\SI{8}{\milli\joule}$.
Determine a tensão em seus terminais.
\resposta{$U\approx\SI{12.6}{\volt}$}
\resolucao{De $W=\frac12CU^{2}$,
\[
U=\sqrt{\frac{2W}{C}}=\sqrt{\frac{2(\pot{8}{-3})}{\pot{100}{-6}}}
=\sqrt{160}=\SI{12.65}{\volt}.
\]
A raiz quadrada indica que a energia é pouco sensível à tensão em valores
baixos e muito sensível em valores altos: dobrar a tensão \emph{quadruplica}
a energia armazenada.}
\end{questao}

\begin{questao}[1]{}
Sobre a tensão em um capacitor ideal, é correto afirmar que ela:
\begin{alternativas}[2]
\item pode variar instantaneamente
\item não pode variar instantaneamente
\item é sempre igual à da fonte
\item é sempre nula em regime permanente CC
\end{alternativas}
\resposta{(b)}
\resolucao{Um salto de tensão implicaria
$\mathrm{d}u/\mathrm{d}t\to\infty$ e, por
$i=C\,\mathrm{d}u/\mathrm{d}t$, corrente infinita --- impossível. Logo
$u(0^{+})=u(0^{-})$.\\
\textbf{Dualidade com o indutor:} no capacitor é a \emph{tensão} que é
contínua; no indutor, a \emph{corrente}. Em regime CC o capacitor tem corrente
nula (comporta-se como circuito aberto) e tensão em geral não nula --- o que
descarta a alternativa (d).}
\end{questao}

\begin{questao}[1]{}
Determine a constante de tempo de um circuito RC com
$R=\SI{4.7}{\kilo\ohm}$ e $C=\SI{10}{\micro\farad}$.
\resposta{$\tau=\SI{47}{\milli\second}$}
\resolucao{
\[
\tau=RC=4700\cdot\pot{10}{-6}=\pot{4.7}{-2}\,\si{\second}
=\SI{47}{\milli\second}.
\]
\textbf{Atenção à estrutura:} no RC, resistência \emph{maior} torna o
transitório \emph{mais lento} ($\tau=RC$); no RL ocorre o contrário
($\tau=L/R$). Verificar a dimensão ajuda:
$\si{\ohm}\times\si{\farad}=\si{\second}$.}
\end{questao}

\blocoaplicacao

\begin{questao}[2]{}
Dois capacitores de $\SI{10}{\micro\farad}$ em série são associados em paralelo
com um de $\SI{15}{\micro\farad}$. Determine $C_{eq}$.
\resposta{$C_{eq}=\SI{20}{\micro\farad}$}
\resolucao{\textbf{Série} (regra inversa, ao contrário dos resistores):
\[
C_s=\frac{10\cdot10}{20}=\SI{5}{\micro\farad}.
\]
\textbf{Paralelo} (soma direta):
\[
C_{eq}=5+15=\SI{20}{\micro\farad}.
\]
Memorize a inversão: capacitores em \emph{série} combinam como resistores em
paralelo, e vice-versa. A razão física é que a série reduz a capacitância
(equivale a afastar as placas), enquanto o paralelo a aumenta (equivale a somar
áreas).}
\end{questao}

\begin{questao}[2]{}
Capacitores de $\SI{1}{\micro\farad}$ e $\SI{2}{\micro\farad}$ estão em série
sob $\SI{30}{\volt}$. Determine $C_{eq}$, a carga e a tensão em cada um.
\resposta{$C_{eq}=\SI{0.667}{\micro\farad}$; $Q=\SI{20}{\micro\coulomb}$;
$U_1=\SI{20}{\volt}$, $U_2=\SI{10}{\volt}$}
\resolucao{
\[
C_{eq}=\frac{1\cdot2}{3}=\SI{0.667}{\micro\farad};
\qquad
Q=C_{eq}U=0{,}667\cdot30=\SI{20}{\micro\coulomb}.
\]
Em série, a \textbf{carga é a mesma} em todos (é a mesma corrente que os
carregou):
\[
U_1=\frac{20}{1}=\SI{20}{\volt},
\qquad
U_2=\frac{20}{2}=\SI{10}{\volt}.
\]
\textbf{Verificação:} $20+10=\SI{30}{\volt}$ \checkmark.\\
A tensão se reparte \emph{inversamente} à capacitância --- o capacitor
\textbf{menor} fica com a maior tensão. É o oposto da série de resistores, e é
a causa de um problema de projeto tratado adiante.}
\end{questao}

\begin{questao}[2]{}
A tensão sobre um capacitor de $\SI{10}{\micro\farad}$ cresce linearmente à
taxa de $\SI{500}{\volt\per\second}$. Determine a corrente.
\resposta{$i=\SI{5}{\milli\ampere}$}
\resolucao{
\[
i=C\frac{\mathrm{d}u}{\mathrm{d}t}=\pot{10}{-6}\cdot500
=\pot{5}{-3}\,\si{\ampere}=\SI{5}{\milli\ampere}.
\]
Enquanto a rampa durar, a corrente é \emph{constante}. Se a tensão parasse de
variar, a corrente cairia a zero instantaneamente --- o capacitor responde à
\textbf{variação} da tensão, não ao seu valor.}
\end{questao}

\begin{questao}[2]{}
Aplica-se a um capacitor de $\SI{2.2}{\micro\farad}$ uma tensão
\textbf{triangular} que sobe $\SI{10}{\volt}$ em $\SI{1}{\milli\second}$ e
desce igualmente no milissegundo seguinte. Descreva e quantifique a corrente.
\resposta{Onda \textbf{quadrada} de $\pm\SI{22}{\milli\ampere}$}
\resolucao{Na subida,
$\mathrm{d}u/\mathrm{d}t=10/10^{-3}=\SI{10000}{\volt\per\second}$:
\[
i=2{,}2\times10^{-6}\cdot10^{4}=\SI{22}{\milli\ampere}.
\]
Na descida a taxa inverte e $i=\SI{-22}{\milli\ampere}$.\\
\textbf{Resultado geral:} o capacitor \emph{deriva} a tensão --- triangular
entra, quadrada sai. O indutor faz o oposto (integra): tensão quadrada nele
produz corrente triangular. É a dualidade dos dois elementos em sua forma mais
visível no osciloscópio.}
\end{questao}

\begin{questao}[2]{}
Um circuito RC com $\tau=\SI{2}{\milli\second}$ é ligado a
$\SI{20}{\volt}$ com o capacitor inicialmente descarregado. Determine
$u_C$ em $t=\tau$, $2\tau$ e $3\tau$.
\resposta{$\SI{12.64}{\volt}$, $\SI{17.29}{\volt}$ e $\SI{19.00}{\volt}$}
\resolucao{Com $u_C=U\left(1-e^{-t/\tau}\right)$ e $U=\SI{20}{\volt}$:
\[
u(\tau)=20(0{,}632)=\SI{12.64}{\volt};
\quad
u(2\tau)=20(0{,}865)=\SI{17.29}{\volt};
\quad
u(3\tau)=20(0{,}950)=\SI{19.00}{\volt}.
\]
As frações $63{,}2\%$, $86{,}5\%$ e $95{,}0\%$ são as mesmas do circuito RL ---
elas dependem só da forma exponencial, não do tipo de elemento.}
\end{questao}

\begin{questao}[2]{}
No mesmo circuito, determine o instante em que a tensão atinge
$\SI{15}{\volt}$.
\resposta{$t\approx\SI{2.77}{\milli\second}$}
\resolucao{
\[
15=20\left(1-e^{-t/\tau}\right)
\;\Longrightarrow\;
e^{-t/\tau}=0{,}25
\;\Longrightarrow\;
t=\tau\ln4=2(1{,}386)=\SI{2.77}{\milli\second}.
\]
\textbf{Atalho útil:} atingir $75\%$ leva exatamente $\ln4\approx1{,}39\tau$;
atingir $50\%$ leva $\ln2\approx0{,}69\tau$. Esse último valor é a "meia-vida"
do transitório e vale a pena memorizar.}
\end{questao}

\begin{questao}[2]{}
Uma fonte de $\SI{24}{\volt}$ alimenta $R_1=\SI{8}{\kilo\ohm}$ em série com
$R_2=\SI{4}{\kilo\ohm}$. Um capacitor está ligado em paralelo com $R_2$.
Determine a tensão no capacitor em regime permanente e a corrente que o
atravessa.
\resposta{$U_C=\SI{8}{\volt}$; $i_C=0$}
\resolucao{Em regime CC o capacitor é um \textbf{circuito aberto}: não conduz.
Toda a corrente passa pelos dois resistores, e a tensão do capacitor é a de
$R_2$:
\[
U_C=24\cdot\frac{4}{8+4}=\SI{8}{\volt};
\qquad
i_C=0.
\]
\textbf{Método:} para determinar o estado final de qualquer circuito CC com
capacitores, basta \emph{removê-los} do esquema e resolver a rede resistiva
restante. Para indutores, o procedimento dual é substituí-los por curtos.}
\end{questao}

\begin{questao}[2]{}
Dois capacitores armazenam a mesma energia. O primeiro tem
$C_1=\SI{100}{\micro\farad}$ sob $\SI{10}{\volt}$. Se
$C_2=\SI{25}{\micro\farad}$, determine sua tensão.
\resposta{$U_2=\SI{20}{\volt}$}
\resolucao{
\[
\frac12C_1U_1^{2}=\frac12C_2U_2^{2}
\;\Longrightarrow\;
U_2=U_1\sqrt{\frac{C_1}{C_2}}=10\sqrt{4}=\SI{20}{\volt}.
\]
Ambos armazenam $\SI{5}{\milli\joule}$. Reduzir $C$ por um fator $4$ exige
\emph{dobrar} a tensão --- e a tensão é justamente o parâmetro caro e limitado
em capacitores reais, pois exige dielétrico mais espesso, o que reduz $C$.}
\end{questao}

\begin{questao}[2]{}
Um capacitor plano tem capacitância $C$. Determine a nova capacitância se:
\begin{itensq}
\item a área das placas for dobrada;
\item a distância entre placas for reduzida à metade;
\item ambas as alterações forem feitas.
\end{itensq}
\resposta{(a) $2C$; (b) $2C$; (c) $4C$}
\resolucao{De $C=\dfrac{\varepsilon A}{d}$: $C\propto A$ e
$C\propto1/d$.\\
(a) $A\to2A$ dá $2C$. (b) $d\to d/2$ dá $2C$. (c) os efeitos se multiplicam:
$4C$.\\
\textbf{Limite prático de (b):} reduzir $d$ aumenta o campo
$E=U/d$ para a mesma tensão. Chega-se rapidamente à rigidez dielétrica do
material, e o capacitor perfura. Há, portanto, um teto para essa estratégia ---
motivo pelo qual capacitores de alta capacitância usam grandes áreas
(folhas enroladas) em vez de dielétricos ultrafinos.}
\end{questao}

\begin{questao}[2]{}
Um capacitor carregado a $\SI{50}{\volt}$ é descarregado através de
$R=\SI{2.2}{\kilo\ohm}$. Determine a corrente inicial e explique sua evolução.
\resposta{$i_0=\SI{22.7}{\milli\ampere}$, decaindo exponencialmente}
\resolucao{No instante inicial a tensão ainda é $\SI{50}{\volt}$
(continuidade), e ela se aplica integralmente ao resistor:
\[
i_0=\frac{50}{2200}=\SI{22.7}{\milli\ampere}.
\]
À medida que o capacitor se descarrega, sua tensão cai e a corrente a acompanha:
$i(t)=i_0e^{-t/\tau}$.\\
\textbf{Observação:} a corrente é máxima no início e a tensão é máxima no
início --- diferentemente do circuito RL, em que a corrente parte de zero. Essa
diferença decorre justamente de qual grandeza é contínua em cada elemento.}
\end{questao}

\blocoaprofundamento

\begin{questao}[3]{}
$C_1=\SI{4}{\micro\farad}$ e $C_2=\SI{12}{\micro\farad}$, ambos com tensão
máxima de trabalho de $\SI{50}{\volt}$, são ligados em série.
\begin{itensq}
\item Determine $C_{eq}$ e a distribuição de tensão sob $\SI{100}{\volt}$.
\item Determine a máxima tensão que a associação suporta.
\item Explique por que o resultado é contraintuitivo.
\end{itensq}
\resposta{(a) $\SI{3}{\micro\farad}$; $\SI{75}{\volt}$ e $\SI{25}{\volt}$;
(b) $\SI{66.7}{\volt}$}
\resolucao{(a) $C_{eq}=\dfrac{4\cdot12}{16}=\SI{3}{\micro\farad}$. Como a carga
é comum, $U\propto1/C$:
\[
U_1=100\cdot\frac{12}{16}=\SI{75}{\volt},
\qquad
U_2=100\cdot\frac{4}{16}=\SI{25}{\volt}.
\]
(b) Sob $\SI{100}{\volt}$, $C_1$ já está a $\SI{75}{\volt}$ --- $50\%$ acima do
limite. Impondo $U_1\le\SI{50}{\volt}$:
\[
U_{\max}=50\cdot\frac{16}{12}=\SI{66.7}{\volt}.
\]
(c) \textbf{É contraintuitivo} porque associar dois componentes de
$\SI{50}{\volt}$ em série sugere $\SI{100}{\volt}$ de capacidade --- e isso
\emph{só} é verdade se eles forem \textbf{iguais}. Com valores diferentes, o
menor capacitor absorve a maior parte da tensão e limita todo o conjunto.\\
\textbf{Regra prática:} em série de capacitores, use sempre unidades idênticas
--- e ainda assim adicione resistores de equalização, como se detalha no bloco
de desafio.}
\end{questao}

\begin{questao}[3]{}
Capacitores de $\SI{2}{\micro\farad}$, $\SI{3}{\micro\farad}$ e
$\SI{5}{\micro\farad}$ estão em paralelo sob $\SI{20}{\volt}$. Determine a
carga de cada um, a carga total e $C_{eq}$, e compare com o caso série.
\resposta{$40$, $60$ e $\SI{100}{\micro\coulomb}$; total
$\SI{200}{\micro\coulomb}$; $C_{eq}=\SI{10}{\micro\farad}$}
\resolucao{Em paralelo a \textbf{tensão} é comum:
\[
Q_k=C_kU \Rightarrow 40,\ 60,\ \SI{100}{\micro\coulomb};
\qquad
Q_{tot}=\SI{200}{\micro\coulomb}.
\]
\[
C_{eq}=2+3+5=\SI{10}{\micro\farad}
\quad\text{(confere: }200/20=10\text{)}.
\]
\textbf{Comparação com a série:} lá, a \emph{carga} é comum e a tensão se
divide; aqui, a \emph{tensão} é comum e a carga se divide. O paralelo soma
capacitâncias e mantém a tensão de trabalho da menor unidade; a série reduz a
capacitância e (idealmente) soma as tensões. São estratégias opostas para
problemas opostos.}
\end{questao}

\begin{questao}[3]{}
Uma fonte de $\SI{12}{\volt}$ alimenta $R_1=\SI{6}{\kilo\ohm}$ em série com
$R_2=\SI{3}{\kilo\ohm}$, e um capacitor de $\SI{100}{\nano\farad}$ está em
paralelo com $R_2$.
\begin{itensq}
\item Determine a tensão final no capacitor.
\item Determine $\tau$ pelo equivalente de Thévenin.
\item Escreva $u_C(t)$.
\end{itensq}
\resposta{(a) $\SI{4}{\volt}$; (b) $\tau=\SI{0.2}{\milli\second}$;
(c) $u_C=4\left(1-e^{-t/0{,}2\,\mathrm{ms}}\right)$}
\resolucao{(a) Com o capacitor aberto em regime,
$U_C=12\cdot\dfrac{3}{9}=\SI{4}{\volt}$.\\
(b) Zerando a fonte, o capacitor vê $R_1\parallel R_2$:
\[
R_{Th}=\frac{6\cdot3}{9}=\SI{2}{\kilo\ohm}
\;\Longrightarrow\;
\tau=R_{Th}C=2000\cdot\pot{100}{-9}=\SI{0.2}{\milli\second}.
\]
(c) $u_C(t)=4\left(1-e^{-t/0{,}0002}\right)$ volts.\\
\textbf{Erro comum:} usar $\tau=R_1C$ ou $\tau=R_2C$. A constante de tempo é
governada pela resistência \emph{vista pelo capacitor} com as fontes zeradas
--- aqui, o paralelo das duas, que é menor que qualquer uma delas. O transitório
é, portanto, mais rápido do que qualquer estimativa ingênua sugeriria.}
\end{questao}

\begin{questao}[3]{}
Um capacitor já está a $\SI{2}{\volt}$ quando é ligado a um circuito que o
levaria a $\SI{10}{\volt}$, com $\tau=\SI{5}{\milli\second}$.
\begin{itensq}
\item Escreva $u_C(t)$.
\item Determine a tensão em $t=\SI{5}{\milli\second}$.
\item Determine a corrente inicial, sabendo que $R=\SI{5}{\kilo\ohm}$.
\end{itensq}
\resposta{(a) $u_C=10-8e^{-t/\tau}$; (b) $\SI{7.06}{\volt}$;
(c) $\SI{1.6}{\milli\ampere}$}
\resolucao{(a) A forma geral de qualquer transitório de primeira ordem é
\[
u(t)=U_f+\left(U_0-U_f\right)e^{-t/\tau}=10-8e^{-t/\tau}.
\]
(b) $u(5\,\mathrm{ms})=10-8e^{-1}=10-2{,}94=\SI{7.06}{\volt}$.\\
(c) Em $t=0^{+}$ a tensão do capacitor é $\SI{2}{\volt}$, e sobre o resistor
sobram $10-2=\SI{8}{\volt}$:
\[
i_0=\frac{8}{5000}=\SI{1.6}{\milli\ampere}.
\]
\textbf{Vantagem da forma geral:} ela cobre carga ($U_0=0$), descarga
($U_f=0$) e qualquer condição intermediária. Basta identificar os três
parâmetros $U_0$, $U_f$ e $\tau$ --- não há três fórmulas a memorizar.}
\end{questao}

\begin{questao}[3]{}
Um capacitor de $\SI{470}{\micro\farad}$ tem resistência de isolação de
$\SI{2}{\mega\ohm}$. Carregado a $\SI{100}{\volt}$ e deixado em circuito
aberto:
\begin{itensq}
\item Determine a constante de tempo de autodescarga.
\item Determine a tensão após $\SI{15}{\minute}$.
\item Comente a implicação de segurança.
\end{itensq}
\resposta{(a) $\tau=\SI{940}{\second}\approx\SI{15.7}{\minute}$;
(b) $\approx\SI{38}{\volt}$}
\resolucao{(a) A fuga é modelada por um resistor em paralelo:
\[
\tau=R_{iso}C=\pot{2}{6}\cdot\pot{470}{-6}=\SI{940}{\second}.
\]
(b) Com $t=\SI{900}{\second}$:
\[
u=100\,e^{-900/940}=100\,e^{-0{,}957}=\SI{38.4}{\volt}.
\]
(c) \textbf{Depois de quinze minutos desligado, o capacitor ainda guarda
$\SI{38}{\volt}$} --- e, em bancos de fonte chaveada com centenas de volts, o
resíduo é letal por horas. É por isso que equipamentos com barramento CC trazem
\emph{resistores de descarga} (bleeders) permanentes e o aviso de aguardar
antes de abrir o gabinete.\\
\textbf{Nota:} capacitores eletrolíticos ainda apresentam \emph{absorção
dielétrica} --- após descarga completa, a tensão "ressurge" parcialmente em
minutos. Curto-circuitar antes de manusear é procedimento obrigatório, não
precaução exagerada.}
\end{questao}

\begin{questao}[3]{}
Um capacitor eletrolítico tem ESR (resistência equivalente em série) de
$\SI{0.1}{\ohm}$ e conduz corrente de ondulação de
$\SI{1.5}{\ampere}$ eficazes.
\begin{itensq}
\item Determine a potência dissipada internamente.
\item Explique por que essa dissipação limita a vida do componente.
\item Indique duas formas de reduzi-la.
\end{itensq}
\resposta{(a) $P=\SI{0.225}{\watt}$}
\resolucao{(a)
\[
P=\mathrm{ESR}\cdot I_{rms}^{2}=0{,}1(1{,}5)^{2}=\SI{0.225}{\watt}.
\]
(b) Essa potência é dissipada \emph{dentro} do capacitor, elevando a
temperatura do eletrólito. A vida útil de eletrolíticos segue uma regra
empírica severa: ela \textbf{cai pela metade a cada $\SI{10}{\celsius}$} de
elevação. Um componente de $\SI{5000}{\hour}$ a $\SI{105}{\celsius}$ dura
$\SI{20000}{\hour}$ a $\SI{85}{\celsius}$ --- e o inverso também vale.
Pior: o eletrólito evapora, a ESR \emph{aumenta}, a dissipação cresce e o
processo se realimenta.\\
(c) \textbf{(i) Dividir a corrente:} usar $n$ capacitores em paralelo reduz a
corrente por unidade a $I/n$ e a dissipação individual a $P/n^{2}$.
Com dois, cada um dissipa $\SI{56}{\milli\watt}$ --- um quarto do original.\\
\textbf{(ii) Escolher tecnologia de baixa ESR:} capacitores de polímero sólido
têm ESR uma ordem de grandeza menor, ao custo de capacitância e tensão
menores.\\
\textbf{Consequência de diagnóstico:} em fontes chaveadas, o capacitor de saída
é quase sempre o primeiro componente a falhar --- e a ESR crescente, não a
perda de capacitância, é o sintoma que primeiro aparece.}
\end{questao}

\begin{questao}[3]{}
Um capacitor de $\SI{100}{\micro\farad}$, inicialmente descarregado, é
carregado a $\SI{20}{\volt}$ através de um resistor.
\begin{itensq}
\item Determine a energia fornecida pela fonte.
\item Determine a energia armazenada no capacitor.
\item Determine a energia dissipada no resistor e comente.
\end{itensq}
\resposta{(a) $\SI{40}{\milli\joule}$; (b) $\SI{20}{\milli\joule}$;
(c) $\SI{20}{\milli\joule}$ --- exatamente metade}
\resolucao{(a) A fonte mantém $\SI{20}{\volt}$ enquanto entrega toda a carga
final $Q=CU=\SI{2}{\milli\coulomb}$:
\[
W_{fonte}=QU=CU^{2}=\pot{100}{-6}(400)=\SI{40}{\milli\joule}.
\]
(b) $W_C=\frac12CU^{2}=\SI{20}{\milli\joule}$.\\
(c) Por conservação, $W_R=40-20=\SI{20}{\milli\joule}$.\\
\textbf{Resultado notável:} o rendimento do carregamento de um capacitor por
fonte de tensão através de resistor é \textbf{sempre $50\%$} --- não importa o
valor de $R$, nem de $C$, nem da tensão. Diminuir $R$ apenas acelera o
processo; a energia perdida é a mesma. A demonstração formal por integração
está no bloco de desafio.\\
\textbf{Consequência prática:} carregadores eficientes não usam resistor em série
--- usam conversores chaveados, que transferem energia em pacotes através de um
indutor e escapam desse limite de $50\%$.}
\end{questao}

\begin{questao}[3]{}
Um capacitor plano usa dielétrico com rigidez de
$\SI{20}{\kilo\volt\per\milli\meter}$ e espessura de
$\SI{0.1}{\milli\meter}$.
\begin{itensq}
\item Determine a tensão de ruptura.
\item Especifique a tensão nominal com margem de $50\%$.
\item Explique por que a rigidez limita o produto capacitância-tensão.
\end{itensq}
\resposta{(a) $\SI{2}{\kilo\volt}$; (b) $\SI{1}{\kilo\volt}$}
\resolucao{(a) $U_{rup}=E_{rig}\cdot d=20\cdot0{,}1=\SI{2}{\kilo\volt}$.\\
(b) Com margem de $50\%$, especifica-se $\SI{1}{\kilo\volt}$ --- e ainda assim
é prudente derrateá-la em alta temperatura e sob ondulação.\\
(c) Combinando $C=\dfrac{\varepsilon A}{d}$ com
$U_{\max}=E_{rig}d$:
\[
C\,U_{\max}=\varepsilon A E_{rig}
\quad\Longrightarrow\quad
W_{\max}=\frac12CU_{\max}^{2}=\frac12\varepsilon E_{rig}^{2}\,(A\,d).
\]
O termo $Ad$ é o \textbf{volume} de dielétrico. Ou seja, a energia máxima
armazenável depende apenas do volume e das propriedades do material
($\varepsilon$ e $E_{rig}$) --- \emph{não} da geometria escolhida.\\
\textbf{Consequência:} não existe truque geométrico para ganhar energia. Só
avança quem melhora o material. É por isso que o desenvolvimento de capacitores
é, essencialmente, ciência de materiais.}
\end{questao}

\begin{questao}[3]{}
Em um circuito RC de carga, mediu-se tensão final de $\SI{10}{\volt}$,
corrente inicial de $\SI{2}{\milli\ampere}$ e tensão de $\SI{6.32}{\volt}$ em
$\SI{3}{\milli\second}$. Determine $R$ e $C$.
\resposta{$R=\SI{5}{\kilo\ohm}$; $C=\SI{0.6}{\micro\farad}$}
\resolucao{\textbf{Resistência:} em $t=0^{+}$ o capacitor está descarregado e
toda a tensão cai em $R$:
\[
R=\frac{U_f}{i_0}=\frac{10}{\pot{2}{-3}}=\SI{5}{\kilo\ohm}.
\]
\textbf{Constante de tempo:} $6{,}32/10=0{,}632$ --- exatamente a fração de
$1\tau$. Logo $\tau=\SI{3}{\milli\second}$ e
\[
C=\frac{\tau}{R}=\frac{\pot{3}{-3}}{5000}=\pot{6}{-7}\,\si{\farad}
=\SI{0.6}{\micro\farad}.
\]
\textbf{Estrutura da medição:} a corrente inicial isola $R$; o transitório
isola $\tau$; e só a combinação entrega $C$. Nenhuma medida sozinha determina
o capacitor --- é o mesmo padrão do circuito RL, em que o regime permanente dá
$R$ e o transitório dá $L$.}
\end{questao}

\begin{questao}[3]{}
Compare o comportamento de capacitor e indutor preenchendo o raciocínio para os
instantes $t=0^{+}$ e $t\to\infty$, e justifique cada equivalência.
\resposta{Capacitor: curto em $t=0^{+}$ (se descarregado), aberto em regime;
indutor: aberto em $t=0^{+}$ (se sem corrente), curto em regime}
\resolucao{\textbf{Capacitor descarregado em $t=0^{+}$:} $u_C=0$, e um elemento
com tensão nula equivale a um \textbf{curto-circuito}. Em regime CC,
$\mathrm{d}u/\mathrm{d}t=0$ implica $i=0$: equivale a \textbf{circuito
aberto}.\\
\textbf{Indutor sem corrente em $t=0^{+}$:} $i_L=0$, o que equivale a
\textbf{circuito aberto}. Em regime, $\mathrm{d}i/\mathrm{d}t=0$ implica
$v=0$: equivale a \textbf{curto}.\\
\textbf{Atenção à condição inicial:} as equivalências em $t=0^{+}$ pressupõem
estado inicial nulo. Se o capacitor já estiver a $U_0$, ele equivale a uma
\emph{fonte de tensão} $U_0$; se o indutor já conduzir $I_0$, equivale a uma
\emph{fonte de corrente} $I_0$.\\
\textbf{Por que isso é útil:} esses quatro substitutos permitem obter os
valores inicial e final de qualquer transitório de primeira ordem resolvendo
apenas circuitos \emph{resistivos}. Junto com $\tau$, eles determinam a solução
completa sem integrar nada.}
\end{questao}

\begin{questao}[3]{}
Um filtro RC passa-baixas tem $R=\SI{10}{\kilo\ohm}$ e
$C=\SI{47}{\nano\farad}$.
\begin{itensq}
\item Determine a frequência de corte.
\item Relacione-a com a constante de tempo.
\item Descreva o comportamento muito abaixo e muito acima de $f_c$.
\end{itensq}
\resposta{(a) $f_c\approx\SI{339}{\hertz}$; (b) $f_c=1/(2\pi\tau)$}
\resolucao{(a)
\[
f_c=\frac{1}{2\pi RC}=\frac{1}{2\pi(10^{4})(\pot{47}{-9})}
=\SI{339}{\hertz}.
\]
(b) Como $\tau=RC=\SI{0.47}{\milli\second}$, vale
$f_c=\dfrac{1}{2\pi\tau}$: um circuito \emph{lento} (τ grande) filtra a partir
de frequências \emph{baixas}. As duas descrições --- temporal e frequencial ---
são a mesma informação.\\
(c) \textbf{Muito abaixo de $f_c$:} o capacitor tem reatância alta, quase não
carrega, e o sinal passa praticamente intacto (ganho $\approx1$).\\
\textbf{Muito acima:} o capacitor não tem tempo de carregar entre as variações;
a saída é fortemente atenuada, a $-\SI{20}{\decibel}$ por década.\\
\textbf{Exatamente em $f_c$:} o ganho vale $1/\sqrt2\approx0{,}707$
($-\SI{3}{\decibel}$), e a defasagem é de $\SI{45}{\degree}$. É a definição
convencional de frequência de corte.}
\end{questao}

\blocodesafio

\begin{questao}[4]{}
\textbf{(Dedução --- resposta do circuito RC.)} Um RC série com fonte $U$ tem a
chave fechada em $t=0$ e $u_C(0)=0$.
\begin{itensq}
\item Escreva a equação diferencial e resolva por separação de variáveis.
\item Obtenha $i(t)$ e verifique a LKT nos dois extremos.
\item Compare a estrutura da solução com a do circuito RL.
\end{itensq}
\resposta{$u_C=U\left(1-e^{-t/RC}\right)$;
$i=\dfrac{U}{R}e^{-t/RC}$}
\resolucao{(a) Pela LKT, com $i=C\,\mathrm{d}u_C/\mathrm{d}t$:
\[
U=Ri+u_C=RC\frac{\mathrm{d}u_C}{\mathrm{d}t}+u_C
\;\Longrightarrow\;
\frac{\mathrm{d}u_C}{U-u_C}=\frac{\mathrm{d}t}{RC}.
\]
Integrando de $0$ a $t$:
\[
-\ln\!\left(\frac{U-u_C}{U}\right)=\frac{t}{RC}
\;\Longrightarrow\;
u_C(t)=U\left(1-e^{-t/RC}\right).
\]
(b)
\[
i=C\frac{\mathrm{d}u_C}{\mathrm{d}t}
=CU\cdot\frac{1}{RC}e^{-t/RC}=\frac{U}{R}e^{-t/RC}.
\]
\emph{Em $t=0^{+}$:} $u_C=0$ e $Ri=U$; soma $=U$ \checkmark.\\
\emph{Em $t\to\infty$:} $u_C=U$ e $Ri=0$; soma $=U$ \checkmark.\\
(c) \textbf{Dualidade estrutural.} No RL, a \emph{corrente} sobe como
$1-e^{-t/\tau}$ e a tensão do indutor decai como $e^{-t/\tau}$. No RC, é a
\emph{tensão} que sobe e a corrente que decai. Trocando
$L\leftrightarrow C$, $i\leftrightarrow u$ e $R\leftrightarrow1/R$, uma solução
vira a outra --- razão pela qual basta dominar um dos dois casos.}
\end{questao}

\begin{questao}[4]{}
\textbf{(Demonstração --- rendimento do carregamento.)} Um capacitor
descarregado é carregado por uma fonte $U$ através de $R$.
\begin{itensq}
\item Integre $p_R=Ri^{2}$ para obter a energia dissipada.
\item Mostre que ela vale $\frac12CU^{2}$, independentemente de $R$.
\item Explique como conversores chaveados contornam esse limite.
\end{itensq}
\resposta{$W_R=\frac12CU^{2}$; rendimento sempre $50\%$}
\resolucao{(a) Com $i=\dfrac{U}{R}e^{-t/\tau}$ e $\tau=RC$:
\[
W_R=\int_0^{\infty}R\left(\frac{U}{R}\right)^{2}e^{-2t/\tau}\,\mathrm{d}t
=\frac{U^{2}}{R}\left[-\frac{\tau}{2}e^{-2t/\tau}\right]_0^{\infty}
=\frac{U^{2}\tau}{2R}.
\]
(b) Substituindo $\tau=RC$:
\[
W_R=\frac{U^{2}}{2R}\cdot RC=\frac12CU^{2}.
\]
O $R$ \textbf{cancela}: ele aparece na potência ($\propto1/R$) e desaparece na
duração ($\tau\propto R$). Como a fonte forneceu $CU^{2}$ e o capacitor reteve
$\frac12CU^{2}$, o rendimento é exatamente $50\%$ --- sempre.\\
(c) \textbf{O limite decorre de forçar a tensão do capacitor a saltar de $0$
para $U$ contra uma fonte rígida.} Conversores chaveados evitam isso: um
indutor recebe energia da fonte em regime de corrente, e depois a entrega ao
capacitor. Como o indutor é um elemento reativo (não dissipa), a transferência
pode, em princípio, ser feita sem perda --- e na prática rendimentos de $90\%$
a $95\%$ são rotineiros.\\
\textbf{Observação:} o mesmo resultado de $50\%$ aparece no compartilhamento de
carga entre dois capacitores e na descarga de um indutor. Em todos os casos, a
independência em relação a $R$ tem a mesma origem: potência e duração variam
inversamente.}
\end{questao}

\begin{questao}[4]{}
\textbf{(Projeto de banco.)} Deseja-se um banco de $\SI{100}{\micro\farad}$
capaz de operar em $\SI{100}{\volt}$, dispondo apenas de capacitores de
$\SI{100}{\micro\farad}$ e $\SI{50}{\volt}$.
\begin{itensq}
\item Determine a associação e o número de unidades.
\item Explique por que são necessários resistores de equalização.
\item Dimensione-os para uma corrente de fuga de até
      $\SI{10}{\micro\ampere}$.
\end{itensq}
\resposta{(a) duas séries de dois, em paralelo --- 4 unidades;
(c) $R\le\SI{50}{\kilo\ohm}$, dissipando $\SI{50}{\milli\watt}$ cada}
\resolucao{(a) \textbf{Série de dois:} suporta $\SI{100}{\volt}$, mas a
capacitância cai para $\SI{50}{\micro\farad}$.\\
\textbf{Duas dessas séries em paralelo:} $50+50=\SI{100}{\micro\farad}$,
mantendo $\SI{100}{\volt}$. Total: \textbf{4 capacitores}.\\
Note o custo: quadruplicou-se o número de componentes para dobrar a tensão. É o
preço inevitável --- a energia armazenada, essa sim, quadruplicou.\\
(b) \textbf{Sem equalização, a divisão de tensão em série é instável.} A tensão
de repouso é fixada pelas \emph{correntes de fuga}, não pelas capacitâncias: o
capacitor de menor fuga assume mais tensão. Como a fuga cresce com a tensão e
com a temperatura, o desequilíbrio se realimenta e pode levar um dos capacitores
a ultrapassar $\SI{50}{\volt}$ e perfurar --- o que joga toda a tensão no outro
e o destrói em seguida.\\
(c) Ligue um resistor em paralelo com cada capacitor, dimensionado para que sua
corrente domine a fuga. Adotando um fator $100$:
\[
I_R=100\cdot\pot{10}{-6}=\SI{1}{\milli\ampere}
\;\Longrightarrow\;
R\le\frac{50}{\pot{1}{-3}}=\SI{50}{\kilo\ohm}.
\]
Dissipação: $P=\dfrac{50^{2}}{\pot{50}{3}}=\SI{50}{\milli\watt}$ por resistor
--- especifique $\frac14\,\si{\watt}$.\\
\textbf{Bônus:} esses resistores também funcionam como \emph{bleeders},
descarregando o banco quando o equipamento é desligado. Com
$\tau=RC=50\,\mathrm{k}\times100\,\mu=\SI{5}{\second}$, o banco fica seguro em
menos de meio minuto.}
\end{questao}

\begin{questao}[4]{}
\textbf{(Projeto de filtro.)} Projete um passa-baixas RC de primeira ordem com
$f_c=\SI{100}{\hertz}$, usando $C=\SI{100}{\nano\farad}$.
\begin{itensq}
\item Determine $R$ e escolha o valor comercial E24.
\item Recalcule $f_c$ e avalie o erro.
\item Determine a atenuação em $\SI{1}{\kilo\hertz}$ e discuta a limitação da
      primeira ordem.
\end{itensq}
\resposta{(a) $R=\SI{15.9}{\kilo\ohm}\to\SI{16}{\kilo\ohm}$;
(b) $f_c=\SI{99.5}{\hertz}$, erro $0{,}5\%$;
(c) $\approx-\SI{20}{\decibel}$}
\resolucao{(a)
\[
R=\frac{1}{2\pi f_cC}=\frac{1}{2\pi(100)(\pot{100}{-9})}
=\SI{15915}{\ohm}.
\]
O valor E24 mais próximo é $\SI{16}{\kilo\ohm}$.\\
(b) $f_c=\dfrac{1}{2\pi(16000)(\pot{100}{-9})}=\SI{99.5}{\hertz}$ --- erro de
$0{,}5\%$, muito abaixo da tolerância do próprio capacitor (tipicamente
$\pm10\%$ ou pior). \textbf{O capacitor, e não o resistor, domina a incerteza
do projeto.}\\
(c) Uma década acima de $f_c$, o ganho é
\[
|H|=\frac{1}{\sqrt{1+(f/f_c)^{2}}}\approx\frac{1}{10}
\;\Longrightarrow\;
-\SI{20}{\decibel}.
\]
\textbf{Limitação:} $\SI{20}{\decibel}$ por década é pouco. Para atenuar
$\SI{60}{\decibel}$ seria preciso $f=\SI{100}{\kilo\hertz}$ --- três décadas
acima. Se a especificação exigir corte mais abrupto, é necessário aumentar a
\emph{ordem} do filtro (cascatear seções ou usar topologia ativa), não
apertar $R$ e $C$.\\
\textbf{Cuidado ao cascatear:} duas seções RC ligadas diretamente
\emph{carregam} uma à outra e o $f_c$ resultante não é o previsto. É preciso
desacoplá-las com um buffer, ou projetar a rede como um todo.}
\end{questao}

\begin{questao}[4]{}
\textbf{(Propagação de incerteza.)} Analise o efeito das tolerâncias de $R$ e
$C$ sobre a constante de tempo.
\begin{itensq}
\item Obtenha a expressão da incerteza relativa de $\tau$.
\item Avalie para $R$ com $\pm1\%$ e $C$ eletrolítico com $\pm20\%$, nos
      critérios de pior caso e estatístico.
\item Conclua sobre onde investir em precisão.
\end{itensq}
\resposta{(a) $\dfrac{\Delta\tau}{\tau}=\dfrac{\Delta R}{R}+\dfrac{\Delta C}{C}$;
(b) $21\%$ (pior caso) e $20{,}0\%$ (estatístico)}
\resolucao{(a) Como $\tau=RC$ é um \emph{produto}, tome o logaritmo e
diferencie:
\[
\ln\tau=\ln R+\ln C
\;\Longrightarrow\;
\frac{\Delta\tau}{\tau}=\frac{\Delta R}{R}+\frac{\Delta C}{C}.
\]
As incertezas \textbf{relativas} se somam --- diferentemente da soma, em que se
somam as absolutas.\\
(b) \emph{Pior caso:} $1\%+20\%=21\%$.\\
\emph{Estatístico:} $\sqrt{1^{2}+20^{2}}=\sqrt{401}=20{,}02\%$.\\
(c) \textbf{Os dois critérios praticamente coincidem} --- e isso é a
informação relevante. Quando uma parcela domina de forma tão desproporcional,
combinar em quadratura não ajuda: o resultado é essencialmente a maior
tolerância isolada.\\
\textbf{Conclusão de projeto:} trocar o resistor de $1\%$ por um de $0{,}1\%$
não muda nada ($20{,}00\%$ contra $20{,}02\%$). Todo o ganho está no capacitor.
Se a aplicação exigir temporização precisa, é preciso abandonar o eletrolítico
--- filme de poliéster oferece $\pm5\%$, e polipropileno chega a $\pm1\%$ com
excelente estabilidade térmica.\\
\textbf{Regra geral:} identifique a maior contribuição \emph{antes} de
especificar componentes. Precisão gasta em parcelas menores é dinheiro
jogado fora.}
\end{questao}

\begin{questao}[4]{}
\textbf{(Dedução --- força entre as placas.)} Considere um capacitor plano de
área $A$ e separação $d$.
\begin{itensq}
\item Obtenha a força entre as placas pelo método da energia, com carga
      constante.
\item Expresse o resultado em função de $U$ e discuta o caso a tensão
      constante.
\item Avalie para $A=\SI{100}{\centi\meter\squared}$,
      $d=\SI{1}{\milli\meter}$ e $U=\SI{1}{\kilo\volt}$.
\end{itensq}
\dados{$\varepsilon_0=\pot{8.85}{-12}\,\si{\farad\per\meter}$}
\resposta{$F=\dfrac{\varepsilon_0AU^{2}}{2d^{2}}\approx\SI{44}{\milli\newton}$,
atrativa}
\resolucao{(a) Com o capacitor \textbf{isolado}, $Q$ é constante e convém usar
\[
W=\frac{Q^{2}}{2C}=\frac{Q^{2}d}{2\varepsilon_0A}.
\]
A força que o sistema exerce é
\[
F=-\frac{\mathrm{d}W}{\mathrm{d}d}=-\frac{Q^{2}}{2\varepsilon_0A},
\]
negativa, isto é, \textbf{atrativa} --- as placas tendem a se aproximar,
reduzindo a energia.\\
(b) Substituindo $Q=CU=\dfrac{\varepsilon_0AU}{d}$:
\[
F=\frac{\varepsilon_0AU^{2}}{2d^{2}}\quad\text{(em módulo)}.
\]
\textbf{Sutileza importante:} a tensão constante, a energia \emph{aumenta} ao
aproximar as placas ($W=\frac12CU^{2}$ com $C$ crescendo), o que sugeriria
força repulsiva. O paradoxo se resolve incluindo a bateria: ela fornece
$U\,\mathrm{d}Q$, o dobro da variação de energia do campo. Feito o balanço
completo, a força é atrativa e de \emph{mesmo módulo} --- como tem de ser, pois
a força é uma grandeza física e não depende do que se escolhe manter
constante.\\
(c)
\[
F=\frac{(\pot{8.85}{-12})(\pot{1}{-2})(10^{6})}{2(10^{-6})}
=\SI{44.3}{\milli\newton}.
\]
\textbf{Ordem de grandeza:} cerca de $\SI{4.5}{\gram}$-força para
$\SI{1}{\kilo\volt}$ --- pequena, mas suficiente para fazer as placas vibrarem
audivelmente sob tensão alternada. É o "zumbido" de capacitores cerâmicos
classe 2, e o princípio de funcionamento de microfones e atuadores
eletrostáticos de MEMS.}
\end{questao}

\begin{questao}[4]{}
\textbf{(Projeto de temporizador.)} Um circuito deve atingir $70\%$ da tensão
final em exatamente $\SI{1}{\second}$.
\begin{itensq}
\item Determine $\tau$.
\item Escolha $C$ e determine $R$, com valores comerciais.
\item Avalie a repetibilidade real do temporizador.
\end{itensq}
\resposta{(a) $\tau=\SI{0.83}{\second}$; (b) $C=\SI{100}{\micro\farad}$ e
$R=\SI{8.2}{\kilo\ohm}$; (c) erro dominado pelo capacitor}
\resolucao{(a) De $0{,}70=1-e^{-t/\tau}$:
\[
e^{-t/\tau}=0{,}30
\;\Longrightarrow\;
\tau=\frac{t}{\ln(1/0{,}30)}=\frac{1}{1{,}204}=\SI{0.83}{\second}.
\]
(b) Com $C=\SI{100}{\micro\farad}$:
$R=\tau/C=0{,}83/\pot{100}{-6}=\SI{8306}{\ohm}$. O valor E24 mais próximo é
$\SI{8.2}{\kilo\ohm}$, que dá $\tau=\SI{0.82}{\second}$ e
\[
t_{70\%}=0{,}82(1{,}204)=\SI{0.987}{\second}\quad(\text{erro de }1{,}3\%).
\]
(c) \textbf{O erro de $1{,}3\%$ do arredondamento é irrelevante} diante da
tolerância do eletrolítico ($\pm20\%$), que se transfere integralmente para
$\tau$: o tempo real ficará entre $0{,}79$ e $\SI{1.19}{\second}$. Some-se a
deriva térmica (vários por cento entre $0$ e $\SI{60}{\celsius}$) e o
envelhecimento.\\
\textbf{Conclusão:} um temporizador RC com eletrolítico serve para "cerca de um
segundo", não para "um segundo". Para precisão, as alternativas são capacitor
de filme (levando a $R$ muito maior e $C$ menor) ou --- o caminho moderno ---
um contador digital com oscilador a cristal, cuja estabilidade é de partes por
milhão.\\
\textbf{Observação de método:} note que o requisito foi dado em $70\%$, e não
em $\tau$. Confundir os dois levaria a um erro de $20\%$ no projeto, pois
$t_{70\%}=1{,}204\,\tau$.}
\end{questao}

\begin{questao}[4]{}
\textbf{(Modelagem --- autodescarga.)} Um capacitor real é modelado por $C$
ideal em paralelo com uma resistência de fuga $R_f$.
\begin{itensq}
\item Deduza a lei de decaimento da tensão em circuito aberto.
\item Defina a constante de tempo de autodescarga e relacione-a com a
      especificação de "corrente de fuga" de catálogo.
\item Avalie para $C=\SI{1000}{\micro\farad}$ e
      $R_f=\SI{100}{\kilo\ohm}$, com tensão inicial de $\SI{300}{\volt}$.
\end{itensq}
\resposta{$u=U_0e^{-t/R_fC}$; $\tau=\SI{100}{\second}$; após
$\SI{5}{\minute}$ restam $\approx\SI{15}{\volt}$}
\resolucao{(a) Em circuito aberto, a única corrente é a de fuga:
\[
C\frac{\mathrm{d}u}{\mathrm{d}t}=-\frac{u}{R_f}
\;\Longrightarrow\;
u(t)=U_0e^{-t/R_fC}.
\]
(b) $\tau_f=R_fC$. Catálogos raramente informam $R_f$; informam a
\emph{corrente de fuga} $I_f$ a uma dada tensão. A conversão é imediata:
$R_f=U/I_f$, e portanto
\[
\tau_f=\frac{U\,C}{I_f}.
\]
Note que $\tau_f$ depende do \emph{produto} $R_fC$, e em eletrolíticos $R_f$
cai quando $C$ cresce --- razão pela qual é a autodescarga de capacitores grandes
não é proporcionalmente mais lenta.\\
(c) $\tau_f=\pot{100}{3}\cdot\pot{1000}{-6}=\SI{100}{\second}$. Após
$\SI{300}{\second}$:
\[
u=300\,e^{-3}=300(0{,}0498)=\SI{14.9}{\volt}.
\]
\textbf{Leitura de segurança:} cinco minutos depois de desligado, o barramento
ainda guarda $\SI{15}{\volt}$ --- e, o que é pior, $\SI{0.11}{\joule}$
concentrados. Aos $\SI{60}{\second}$ ainda havia $\SI{110}{\volt}$, plenamente
letais.\\
\textbf{Por isso o resistor de descarga é obrigatório}, e não opcional: com um
bleeder de $\SI{10}{\kilo\ohm}$, $\tau$ cai para $\SI{10}{\second}$ e o
equipamento fica seguro em menos de um minuto, ao custo de
$\SI{9}{\watt}$ apenas enquanto energizado.}
\end{questao}

\begin{questao}[4]{}
\textbf{(ESR e ondulação.)} Um capacitor de filtro de fonte chaveada tem
$\mathrm{ESR}=\SI{50}{\milli\ohm}$ e deve conduzir ondulação de
$\SI{2}{\ampere}$ eficazes.
\begin{itensq}
\item Determine a dissipação e a ondulação de tensão devida \emph{apenas} à
      ESR.
\item Compare com a ondulação devida à capacitância, para
      $C=\SI{470}{\micro\farad}$ a $\SI{100}{\kilo\hertz}$.
\item Conclua sobre o critério de seleção.
\end{itensq}
\resposta{(a) $\SI{0.2}{\watt}$ e $\approx\SI{100}{\milli\volt}$ de pico;
(b) $\approx\SI{5}{\milli\volt}$ --- a ESR domina}
\resolucao{(a) $P=\mathrm{ESR}\cdot I_{rms}^{2}=0{,}05(4)=\SI{0.2}{\watt}$.\\
A ondulação de tensão que a ESR produz acompanha a corrente instantaneamente:
\[
\Delta u_{ESR}\approx\mathrm{ESR}\cdot I_{pico}\approx0{,}05(2)=\SI{100}{\milli\volt}.
\]
(b) A parcela capacitiva depende da carga transferida por semiciclo. Tomando
$\Delta t\approx1/(2f)=\SI{5}{\micro\second}$:
\[
\Delta u_C\approx\frac{I\,\Delta t}{C}
=\frac{2(\pot{5}{-6})}{\pot{470}{-6}}\approx\SI{21}{\milli\volt}.
\]
(c) \textbf{A ESR domina} --- ela responde por cerca de $80\%$ da ondulação
total. Aumentar $C$ para $\SI{1000}{\micro\farad}$ reduziria $\Delta u_C$ pela
metade e deixaria a ondulação praticamente inalterada.\\
\textbf{Critério correto de seleção:} em fontes chaveadas de alta frequência,
escolhe-se o capacitor pela \textbf{ESR e pela corrente de ondulação
admissível}, não pela capacitância. É comum usar vários capacitores menores em
paralelo --- não para somar $C$, mas para dividir a ESR por $n$ e a dissipação
individual por $n^{2}$.\\
\textbf{Nota:} em frequências baixas ($\SI{120}{\hertz}$ de retificação de
rede), a relação se inverte e a capacitância volta a ser o critério dominante.
O parâmetro decisivo depende da faixa de operação.}
\end{questao}

\begin{questao}[4]{}
\textbf{(Comparação tecnológica.)} Um supercapacitor de $\SI{3000}{\farad}$ e
$\SI{2.7}{\volt}$ tem massa de $\SI{0.50}{\kilogram}$.
\begin{itensq}
\item Determine a energia armazenada e a densidade de energia.
\item Compare com uma bateria de íon-lítio ($\approx\SI{250}{\watt\hour\per\kilogram}$).
\item Aponte em que aspecto o supercapacitor é superior e por quê.
\end{itensq}
\resposta{(a) $\SI{10.9}{\kilo\joule}=\SI{3.04}{\watt\hour}$, ou
$\SI{6.1}{\watt\hour\per\kilogram}$; (b) cerca de $40$ vezes menos}
\resolucao{(a)
\[
W=\frac12CU^{2}=\frac12(3000)(2{,}7)^{2}=\SI{10935}{\joule}
=\frac{10935}{3600}=\SI{3.04}{\watt\hour}.
\]
Densidade: $3{,}04/0{,}50=\SI{6.1}{\watt\hour\per\kilogram}$.\\
(b) A bateria armazena cerca de \textbf{40 vezes mais} energia por quilograma.
Para igualar uma célula de íon-lítio de $\SI{20}{\gram}$, seriam necessários
$\SI{0.8}{\kilogram}$ de supercapacitores.\\
(c) \textbf{Potência específica e vida útil.} O supercapacitor pode ser
carregado e descarregado em \emph{segundos}, com correntes de centenas de
ampères, porque armazena energia \textbf{eletrostaticamente} --- sem reações
químicas, sem difusão iônica lenta, sem degradação de eletrodos. Suporta
centenas de milhares de ciclos, contra alguns milhares da bateria, e funciona
bem a baixas temperaturas.\\
\textbf{Consequência de aplicação:} baterias são \emph{reservatórios} de
energia; supercapacitores são \emph{amortecedores} de potência. Por isso
convivem em veículos elétricos --- o supercapacitor absorve o pico da frenagem
regenerativa e entrega o pico da aceleração, poupando a bateria justamente dos
regimes que mais a degradam.\\
\textbf{Ressalva:} a tensão do supercapacitor cai linearmente com a carga
(pois $Q=CU$), ao contrário do platô da bateria. Aproveitar mais de $75\%$ da
energia exige um conversor CC-CC de entrada ampla.}
\end{questao}
